\subsection{Structural parsing with Docling}\label{subsec:doc-parsing-with-docling}

Docling~\parencite{} (Auer et al. 2024) is a visual layout analysis system for PDFs, and is the backbone of the document extraction pipeline. Docling takes a PDF as input, and outputs a typed list of elements: paragraphs, headings, captions, list items, figures, tables, headers, and footers. Each element is annotated with its page, type label, and a bounding box. 

The pipeline uses this output to reconstruct the basic structure of the PDF document. Heading-like elements and their nesting levels are used to create a section hierarchy, while paragraph-like elements become the main textual evidence units consumed downstream. Figure and table bounding boxes provide the initial candidates for extraction, caption linking, and text-to-artifact references. 

Docling's output is the starting point of this pipeline, not its final result. Docling detects tables directly from the page layout, but tables without a recoverable cell grid are outside what a layout parser can reliably recover. The pipeline therefore combines Docling with Microsoft's Table Transformer (TATR), an image-based detector that localizes tables from rendered page images~\parencite{}. A hybrid table detector merges the two sets of table outputs. This detector is the configuration carried into the final pipeline after comparison with Docling-only and TATR-only baselines, the results of which are discussed in Chapter~\ref{chapter:Results}. Docling remains the sole source of body text, section structure, and figure regions -- the hybrid detector is only responsible for table extraction.

The remaining stages refine Docling's text output, rather than replace it. The two-pass ghost text detector of Subsection~\ref{subsec:doc-two-pass} compares Docling's text blocks to an additional, pixel-rendered version of the same region, after dropping texts that are reported by the parser but are not rendered on the page. Docling then parses the redacted page again. This re-extracted layout, rather than a modified version of the first-pass output, is assembled into hierarchical text (Subsection~\ref{subsec:doc-text-assembly}). Unnecessary Docling elements are removed, paragraphs split across columns or pages are merged, and citation markers are stripped. The resulting elements are written to the database, where each element is traceable to its source (Subsection~\ref{subsec:doc-provenance-preserving-storage}).